\section{有向观测：从序列到网与滤子}\label{sec:01-01}

\subsection{有向集、最终性与网}\label{subsec:1-1-1}

序列把逼近过程安放在整数时钟上。这个安排适合逐项部分和，却不适合同时加严许多彼此独立的有限条件。
例如，一次观测也许要求函数在有限个点上接近目标；下一次观测也许换了另一批点。真正需要的“后一步”
不是把某个整数加一，而是把两批点合并。分割的共同加细、有限测试函数的并集和无序和的有限部分和都具有
相同结构。

\begin{example}[有限条件的有向结构]\label{ex:1-1-1-1}
设 \(S\) 为集合，以 \(\Fin(S)\) 表示 \(S\) 的有限子集全体，并按包含关系排序。给定
\(F,G\in\Fin(S)\)，\(F\cup G\) 同时包含二者。若 \((v_s)_{s\in S}\) 位于 Hausdorff 拓扑向量空间，令
\[
  s_F=\sum_{s\in F}v_s,\qquad F\in\Fin(S).
\]
这里每个和都是有限和。若这些有限和在 \(F\) 增大时收敛，则极限没有使用 \(S\) 的编号；这正是无条件
求和所需的语法。拓扑向量空间在本例中只用到有限加法的连续性和极限唯一性。
\end{example}

\begin{definition}[有向集]\label{def:1-1-1-1}
非空集合 \(I\) 上的关系 \(\leq\) 若自反且传递，称为\emph{预序}。若进一步对任意
\(\alpha,\beta\in I\)，存在 \(\gamma\in I\) 使
\[
  \alpha\leq\gamma,\qquad \beta\leq\gamma,
\]
则称 \((I,\leq)\) 为\emph{有向集}。有向关系不要求反对称。
\end{definition}

\begin{definition}[共尾子集]\label{def:1-1-1-cofinal}
有向集 \(I\) 的子集 \(J\) 称为\emph{共尾的}，若对每个 \(\alpha\in I\)，存在 \(\beta\in J\)
满足 \(\alpha\leq\beta\)。这个条件只保证 \(J\) 能越过每个固定阈值；它与定义子网时使用的“最终共尾”
不同。
\end{definition}

\begin{definition}[网、最终与频繁]\label{def:1-1-1-2}
从有向集 \(I\) 到集合 \(X\) 的映射 \(\alpha\mapsto x_\alpha\) 称为 \(X\) 中的\emph{网}。
性质 \(P(x_\alpha)\) \emph{最终成立}，是指存在 \(\alpha_0\in I\)，使每个
\(\alpha\geq\alpha_0\) 都满足 \(P(x_\alpha)\)。性质 \(P(x_\alpha)\) \emph{频繁成立}，是指对每个
\(\alpha_0\in I\)，都存在 \(\alpha\geq\alpha_0\) 满足 \(P(x_\alpha)\)。
\end{definition}

\begin{definition}[网收敛]\label{def:1-1-1-3}
设 \(X\) 为拓扑空间。网 \((x_\alpha)\) \emph{收敛}到 \(x\in X\)，若对 \(x\) 的每个邻域 \(U\)，
关系 \(x_\alpha\in U\) 最终成立。记作 \(x_\alpha\to x\)。
\end{definition}

\begin{definition}[网的聚点]\label{def:1-1-1-cluster}
点 \(x\in X\) 称为网 \((x_\alpha)_{\alpha\in I}\) 的\emph{聚点}，若对 \(x\) 的每个邻域 \(U\)，
关系 \(x_\alpha\in U\) 频繁成立。等价地，对每个 \(\alpha_0\in I\) 和每个
\(U\in\mathcal N(x)\)，
\[
  U\cap\{x_\alpha:\alpha\geq\alpha_0\}\ne\varnothing.
\]
\end{definition}

邻域族本身给出一个重要有向集。规定 \(U\preceq V\) 当且仅当 \(U\supset V\)，指标越大便表示邻域越小，
而 \(U\cap V\) 是共同后继。需要注意，邻域 \(U\) 是 \(X\) 的子集，不是 \(X\) 中的点；从闭包构造网时，
我们将把邻域与其中的点同时放入指标。

\begin{remark}[从序列到 Moore--Smith 网]
十九世纪分析把极限写成序列极限，在度量空间中十分成功。Moore--Smith 网所增加的不是另一套数值计算，
而是把邻域缩小、分割加细和有限条件增多放到同一种“最终同时满足”的语法中。下面的闭包与连续性定理说明，
这种语法恰好恢复了一般拓扑所需的全部极限信息。
\end{remark}

\begin{proposition}[最终性质的有限交法则]\label{proposition:1-1-1-1}
若 \(P(x_\alpha)\) 与 \(Q(x_\alpha)\) 都最终成立，则
\(P(x_\alpha)\wedge Q(x_\alpha)\) 最终成立。同一结论对任意有限多个最终性质成立。
\end{proposition}
\begin{proof}
取阈值 \(\alpha_P,\alpha_Q\)，分别保证 \(P,Q\) 在对应阈值以后成立。由有向性，存在 \(\gamma\) 满足
\(\gamma\geq\alpha_P,\alpha_Q\)。若 \(\alpha\geq\gamma\)，传递性给出
\(\alpha\geq\alpha_P,\alpha_Q\)，故 \(P(x_\alpha)\) 与 \(Q(x_\alpha)\) 同时成立。有限情形由归纳法得到。
\end{proof}

\begin{theorem}[网刻画闭包]\label{theorem:1-1-1-2}
设 \(A\subset X\)，\(x\in X\)。则
\[
  x\in\overline A
  \quad\Longleftrightarrow\quad
  \text{存在取值于 \(A\) 的网收敛到 \(x\)}.
\]
\end{theorem}
\begin{proof}
先设 \(x\in\overline A\)。令
\[
  J=\{(U,a):U\in\mathcal N(x),\ a\in A\cap U\}.
\]
闭包条件保证 \(A\cap U\ne\varnothing\)，所以 \(J\ne\varnothing\)。规定
\((U,a)\preceq(V,b)\) 当且仅当 \(V\subset U\)。给定 \((U,a),(V,b)\)，集合 \(U\cap V\)
是 \(x\) 的邻域，故存在 \(c\in A\cap U\cap V\)，于是 \((U\cap V,c)\) 是共同后继。故 \(J\) 有向。

令 \(z_{(U,a)}=a\)。给定 \(x\) 的邻域 \(W\)，取 \(a_0\in A\cap W\)。当
\((U,a)\succeq(W,a_0)\) 时 \(U\subset W\)，从而 \(z_{(U,a)}=a\in W\)。因此 \(z_j\to x\)。
这里指标集包含所有可行的邻域--点对，没有为整个邻域族预先指定选择函数。

反过来，设 \(A\) 中的网 \(x_\alpha\to x\)。每个 \(x\) 的邻域 \(U\) 最终包含 \(x_\alpha\)，因而
\(A\cap U\ne\varnothing\)。这正是 \(x\in\overline A\)。
\end{proof}

\begin{theorem}[连续性的网判别]\label{theorem:1-1-1-3}
映射 \(f:X\to Y\) 在 \(x\in X\) 连续，当且仅当对每个 \(x_\alpha\to x\)，都有
\(f(x_\alpha)\to f(x)\)。
\end{theorem}
\begin{proof}
若 \(f\) 在 \(x\) 连续，给定 \(f(x)\) 的邻域 \(V\)，存在 \(x\) 的邻域 \(U\) 满足
\(f(U)\subset V\)。由 \(x_\alpha\to x\)，关系 \(x_\alpha\in U\) 最终成立，故
\(f(x_\alpha)\in V\) 最终成立。

反过来，若 \(f\) 在 \(x\) 不连续，则存在 \(f(x)\) 的邻域 \(V\)，使每个 \(x\) 的邻域 \(U\)
都含有某个 \(a\in U\) 满足 \(f(a)\notin V\)。令
\[
  J=\{(U,a):U\in\mathcal N(x),\ a\in U,\ f(a)\notin V\},
\]
并仍按邻域缩小排序。与闭包定理的证明相同，\(J\) 非空且有向，网 \(z_{(U,a)}=a\) 收敛到 \(x\)。
但 \(f(z_j)\notin V\) 对每个 \(j\) 成立，故像网不收敛到 \(f(x)\)，与假设矛盾。
\end{proof}

\begin{proposition}[Hausdorff 判别]\label{proposition:1-1-1-4}
拓扑空间 \(X\) 为 Hausdorff，当且仅当每个网至多有一个极限。
\end{proposition}
\begin{proof}
若 \(X\) 为 Hausdorff，设 \(x_\alpha\to x\) 且 \(x_\alpha\to y\)，其中 \(x\ne y\)。取不交邻域
\(U\ni x,V\ni y\)。网最终进入 \(U\)，也最终进入 \(V\)；由有限交法则，它最终进入
\(U\cap V=\varnothing\)，这是不可能的。

若 \(X\) 不是 Hausdorff，则存在 \(x\ne y\)，使任意邻域 \(U\ni x,V\ni y\) 都相交。令
\[
 J=\{(U,V,z):U\in\mathcal N(x),\ V\in\mathcal N(y),\ z\in U\cap V\}.
\]
按两个邻域同时缩小排序。两个指标的共同后继可从
\((U_1\cap U_2)\cap(V_1\cap V_2)\ne\varnothing\) 中取点得到。令网值为第三坐标。给定 \(x\) 的邻域
\(U_0\)，在指标超过任一以 \(U_0\) 为第一邻域坐标的指标后，网值均落在 \(U_0\)；故该网趋于 \(x\)。
同理它趋于 \(y\)。构造使用全部可行三元组，不需要为所有邻域对预选交点。
\end{proof}

\begin{definition}[第一可数]\label{def:1-1-1-first-countable}
若拓扑空间 \(X\) 的每个点 \(x\) 都有可数邻域基，即存在邻域列 \((V_n)\)，使每个 \(x\) 的邻域都包含
某个 \(V_n\)，则称 \(X\) \emph{第一可数}。
\end{definition}

\begin{proposition}[第一可数空间的序列检验]\label{proposition:1-1-1-first-countable-sequential-tests}
设 \(X\) 第一可数。
\begin{enumerate}
  \item \(x\in\overline A\) 当且仅当存在 \(A\) 中序列收敛到 \(x\)；
  \item 对任意拓扑空间 \(Y\)，映射 \(f:X\to Y\) 在 \(x\) 连续，当且仅当它保持所有趋于 \(x\) 的
  序列之极限。
\end{enumerate}
这并不表示 \(X\) 中每个网都有由 \(\N\) 参数化的子网。
\end{proposition}
\begin{proof}
取 \(x\) 的可数邻域基 \((V_n)\)，并令 \(U_n=V_1\cap\cdots\cap V_n\)。则
\(U_{n+1}\subset U_n\)，且 \((U_n)\) 仍是邻域基。

若 \(x\in\overline A\)，对每个 \(n\) 取 \(a_n\in A\cap U_n\)。给定邻域 \(U\ni x\)，存在 \(N\)
使 \(U_N\subset U\)，于是 \(n\ge N\) 时 \(a_n\in U_n\subset U_N\subset U\)。故 \(a_n\to x\)。
逆向由闭包定义得到。这次逐 \(n\) 选点使用可数选择；若给定邻域基同时给定了选点规则，则不另需选择。

连续映射保持序列极限。反过来，若 \(f\) 在 \(x\) 不连续，取 \(f(x)\) 的邻域 \(V\)，使每个 \(U_n\)
都含有 \(a_n\) 且 \(f(a_n)\notin V\)。上段证明给出 \(a_n\to x\)，而 \(f(a_n)\) 不趋于 \(f(x)\)，矛盾。

最后说明结尾的区别。令 \(I=\Fin(\mathcal P(\N))\)，按包含排序。这里先把稍后
\cref{def:1-1-3-1} 才正式命名的条件写全：映射 \(\phi:\N\to I\) 最终共尾，是指对每个
\(F_0\in I\)，存在 \(N\) 使
\[
  n\ge N\quad\Longrightarrow\quad F_0\subset\phi(n).
\]
若这样的 \(\phi\) 存在，则对每个 \(A\subset\N\)，把 \(F_0=\{A\}\) 代入便有
\(A\in\phi(n)\)（\(n\ge N\)），所以
\[
  \mathcal P(\N)\subset C:=\bigcup_{n\ge1}\phi(n).
\]
对不同的 \(A,B\subset\N\)，规定
\[
  A\prec B
  \quad\Longleftrightarrow\quad
  \min(A\mathbin\triangle B)\in B.
\]
最小差异位置恰属于二者之一，所以任意两个不同集合中恰有一个在前。还需核对传递性：若
\(A\prec B\prec C\)，令
\[
 r=\min(A\mathbin\triangle B),\qquad
 s=\min(B\mathbin\triangle C).
\]
若 \(r=s\)，则 \(A\prec B\) 给出 \(r\in B\)，而 \(B\prec C\) 给出 \(s\notin B\)，矛盾。
若 \(r<s\)，则 \(A,C\) 在所有小于 \(r\) 的位置相同，并且 \(r\notin A\)、\(r\in B=C\)，故
\(A\prec C\)；若 \(s<r\)，同理有 \(s\notin A=B\)、\(s\in C\)，仍得 \(A\prec C\)。因此
\(\prec\) 是 \(\mathcal P(\N)\) 上的严格词典全序。将每个有限集 \(\phi(n)\) 依此序写成
\(A_{n,1},\ldots,A_{n,k_n}\)，再用固定双射 \(\N\leftrightarrow\N^2\)（越界位置填入空集），得到
序列 \((E_j)\) 枚举 \(C\)。令
\[
  B=\{j\in\N:j\notin E_j\}.
\]
则对每个 \(j\) 都有 \(B\ne E_j\)，故 \(B\notin C\)，这与
\(B\in\mathcal P(\N)\subset C\) 矛盾。故这个有向集没有可数最终共尾参数映射；即使取值网是常值网，
也不存在按\cref{def:1-1-3-1}由 \(\N\) 参数化的子网。
\end{proof}

\begin{example}[分割细化]\label{ex:1-1-1-2}
区间上的带标记分割 \((P,\xi)\) 可按底层分割的细化排序。给定两个分割，把两组分点取并，再在每个新区间
内取标签，便得到共同加细。Riemann 和的稳定性要求：存在一个阈值分割，使所有后续细化和所有标签给出的和
都落在规定误差内。本例只说明索引方式；Stieltjes 和是否收敛还需要对被积函数和积分子另加条件。
\end{example}

\begin{example}[闭包中没有序列逼近的点]\label{ex:1-1-1-3}
令 \(I=\mathcal P(\N)\)，\(X=\{0,1\}^{I}\)。暂时给 \(X\) 配置“一个基本邻域只限制有限多个坐标”的
拓扑；这正是 \S\ref{subsec:1-2-2} 将正式定义的乘积拓扑。令 \(A\) 为有限支撑的 \(0/1\) 函数全体，
令 \(\mathbf 1\) 为全一函数。每个只限制有限坐标集 \(F\subset I\) 的 \(\mathbf1\) 邻域都含有
\(\mathbf1_F\in A\)，故 \(\mathbf1\in\overline A\)。网
\[
  F\longmapsto\mathbf1_F,\qquad F\in\Fin(I),
\]
收敛到 \(\mathbf1\)。

另一方面，任取 \(A\) 中序列 \((\mathbf1_{F_n})\)。用上一证明中的规范词典序逐项排列每个有限 \(F_n\)，
再用 \(\N^2\leftrightarrow\N\) 把所得双重有限表写成序列 \((E_j)_{j\ge1}\)；未使用的位置填入
\(\varnothing\)。于是每个 \(F_n\) 的每个元素都在 \((E_j)\) 中出现。令
\[
  i=\{j\in\N:j\notin E_j\}\in I=\mathcal P(\N).
\]
对每个 \(j\) 都有 \(i\ne E_j\)，故 \(i\notin\bigcup_nF_n\)。于是对所有 \(n\)，
\[
  \mathbf1_{F_n}(i)=0\ne1=\mathbf1(i).
\]
集合
\[
  W_i=\{x\in X:x(i)=1\}
\]
只限制坐标 \(i\)，因而按本例给出的基本邻域规则是 \(\mathbf1\) 的开邻域。若
\(\mathbf1_{F_n}\to\mathbf1\)，则应存在 \(N\) 使 \(n\ge N\) 时
\(\mathbf1_{F_n}\in W_i\)，即 \(\mathbf1_{F_n}(i)=1\)；这与上式对每个 \(n\) 的等式矛盾。
因此不必预先调用尚未正式定义的坐标投影连续性，也已证明该序列不收敛到 \(\mathbf1\)。
\end{example}

\begin{figure}[htbp]
\centering
\begin{tikzpicture}[node distance=1.35cm and 1.5cm,>=Stealth,every node/.style={align=center}]
 \node[book node] (n) {邻域缩小\\\(U\cap V\)};
 \node[book node,right=of n] (f) {有限集增大\\\(F\cup G\)};
 \node[book node,right=of f] (p) {分割细化\\公共加细};
 \node[book node,below=of f] (d) {有限阈值有共同后继};
 \node[book node,below=of d] (l) {最终同时满足\\网极限};
 \draw[book arrow] (n)--(d); \draw[book arrow] (f)--(d); \draw[book arrow] (p)--(d);
 \draw[book arrow] (d)--(l);
\end{tikzpicture}
\caption{三个有向系统共享同一个“共同加严”机制。}\label{fig:1-1-1-directed-systems}
\end{figure}

有向关系不必反对称；把预序擅自改成偏序只会加入未使用的条件。“频繁”也不能解释为“无穷多次”，因为
一般指标集没有可数的次数概念。网恢复了序列在一般拓扑中的闭包和连续性角色，却不会把只对可数并成立的
结论推广到任意并。下一节将用网讨论乘积紧性，随后再把同一语言用于函数族极限。

\begin{exercise}
证明有限多个最终性质的合取最终成立。再对序列取 \(P_k(n)\Longleftrightarrow n\ge k\)，说明每个
\(P_k\) 最终成立而 \(\bigwedge_kP_k\) 从不成立。
\end{exercise}
\begin{exercise}
在积有向集上写出“两个变量同时趋近”的定义，并给出两个迭代极限都存在而同时极限不存在的函数例。
\end{exercise}
\begin{exercise}
设 \((v_s)_{s\in S}\) 位于 Banach 空间。证明有限和网收敛，当且仅当对每个 \(\varepsilon>0\)，存在
有限 \(F_0\subset S\)，使每个有限 \(E\subset S\setminus F_0\) 都满足
\(\norm{\sum_{s\in E}v_s}<\varepsilon\)。充分性证明应先对尺度 \(2^{-n}\) 取见证并作有限并，得到递增
\(F_n\)，用序列完备性处理 \(\sum_{s\in F_n}v_s\)，再由同一尾估计回证整个网收敛。
\end{exercise}
\begin{exercise}
证明若同一集合上的两个拓扑拥有完全相同的收敛网，则两个拓扑相同。
\end{exercise}

\subsection{滤子：把尾部行为从索引中分离}\label{subsec:1-1-2}

一个网携带两类信息：参数如何排列，以及哪些集合最终包含网值。很多论证只使用后一类信息。滤子把“最终
落入哪些集合”直接记录成集合族；重参数化造成的冗余随之消失。

\begin{remark}
滤子由 Cartan 学派系统化。网适合写出具体近似，滤子适合合并集合信息；从网到尾滤子会忘掉参数化，
而从滤子返回网若采用成对索引，则无需为每个滤子成员预选一个点。
\end{remark}

\begin{example}[重参数化后尾部信息不变]\label{ex:1-1-2-reparam}
令 \(x_n=1/n\)，并在 \(\N^2\) 的逐坐标次序上令 \(y_{(n,m)}=1/n\)。从 \((n_0,m_0)\) 开始，\(y\)
所能取到的值为 \(\{1/n:n\ge n_0\}\)，与 \(x\) 从 \(n_0\) 开始的尾值相同。第二个时钟没有增加任何
最终信息。
\end{example}

\begin{definition}[滤子]\label{def:1-1-2-1}
集合 \(X\) 上的\emph{滤子}是非空集合族 \(\mathcal F\subset2^X\)，满足
\begin{enumerate}
 \item \(\varnothing\notin\mathcal F\)；
 \item \(A,B\in\mathcal F\) 蕴含 \(A\cap B\in\mathcal F\)；
 \item \(A\in\mathcal F\) 且 \(A\subset B\subset X\) 蕴含 \(B\in\mathcal F\)。
\end{enumerate}
若 \(\mathcal G\supset\mathcal F\)，称 \(\mathcal G\) 比 \(\mathcal F\) 更细。
\end{definition}

\begin{definition}[滤子基]\label{def:1-1-2-2}
非空集合族 \(\mathcal B\subset2^X\) 称为\emph{滤子基}，若
\(\varnothing\notin\mathcal B\)，且对任意 \(B_1,B_2\in\mathcal B\)，存在
\(B_3\in\mathcal B\) 使 \(B_3\subset B_1\cap B_2\)。它生成的滤子为
\[
  \langle\mathcal B\rangle
  =\{A\subset X:\text{存在 }B\in\mathcal B\text{ 使 }B\subset A\}.
\]
\end{definition}

\begin{definition}[滤子收敛与聚点]\label{def:1-1-2-3}
滤子 \(\mathcal F\) \emph{收敛}到 \(x\)，若
\(\mathcal N(x)\subset\mathcal F\)。点 \(x\) 是 \(\mathcal F\) 的\emph{聚点}，若每个
\(F\in\mathcal F\) 与每个 \(U\in\mathcal N(x)\) 都相交。
\end{definition}

这里的 $\mathcal N(x)$ 指所有包含 $x$ 的某个开邻域的集合，而不只指开集本身；它确实是真滤子。
它不含空集，两个成员分别包含的开邻域之交仍是 $x$ 的开邻域，并且“包含一个开邻域”在扩大集合时
保持。因此后面可以把 $\mathcal N(x)$ 代入关于滤子的命题。

\begin{proposition}[尾滤子]\label{proposition:1-1-2-1}
网 \((x_\alpha)_{\alpha\in I}\) 的尾值集合
\(T_\alpha=\{x_\beta:\beta\ge\alpha\}\) 构成滤子基。其生成滤子称为该网的\emph{尾滤子}。
网收敛到 \(x\) 当且仅当尾滤子收敛到 \(x\)；\(x\) 是网的聚点当且仅当它是尾滤子的聚点。
\end{proposition}
\begin{proof}
每个 \(T_\alpha\) 非空。给定 \(T_\alpha,T_\beta\)，取共同后继 \(\gamma\)，则
\(T_\gamma\subset T_\alpha\cap T_\beta\)，故这些尾集是滤子基。

网收敛到 \(x\)，等价于每个 \(U\in\mathcal N(x)\) 包含某个尾集 \(T_\alpha\)，又等价于 \(U\)
属于尾滤子。聚点方面，\(U\) 与每个 \(T_\alpha\) 相交，等价于 \(x_\alpha\in U\) 频繁成立；尾滤子的
任一成员包含某个 \(T_\alpha\)，所以这又等价于 \(U\) 与每个尾滤子成员相交。
\end{proof}

\begin{theorem}[滤子可由网实现]\label{theorem:1-1-2-2}
对 \(X\) 上任一滤子 \(\mathcal F\)，存在 \(X\) 中的网，其尾滤子恰等于 \(\mathcal F\)。构造不需要
为每个滤子成员预选代表点。
\end{theorem}
\begin{proof}
令
\[
 I=\{(F,x):F\in\mathcal F,\ x\in F\},
 \qquad (F,x)\preceq(G,y)\Longleftrightarrow G\subset F.
\]
滤子不含空集，故 \(I\ne\varnothing\)。给定 \((F,x),(G,y)\)，集合
\(F\cap G\in\mathcal F\) 且非空；取 \(z\in F\cap G\)，则 \((F\cap G,z)\) 是共同后继。令
\(z_{(F,x)}=x\)。从任一指标 \((F,x)\) 开始，尾网值全在 \(F\) 中；反过来，每个 \(y\in F\) 都是指标
\((F,y)\) 的网值，且 \((F,x)\preceq(F,y)\)。因此
\[
  \{z_{(G,y)}:(G,y)\succeq(F,x)\}=F.
\]
尾值集合恰遍历 \(\mathcal F\)，故它们生成的向上闭包为
\[
  \{A\subset X:\exists F\in\mathcal F,\ F\subset A\}=\mathcal F.
\]
\end{proof}

\begin{definition}[像滤子]\label{def:1-1-2-image}
若 \(f:X\to Y\)，\(\mathcal F\) 是 \(X\) 上滤子，定义
\[
 f_*\mathcal F=\{B\subset Y:f^{-1}(B)\in\mathcal F\}.
\]
\end{definition}

\begin{proposition}[像滤子的基本性质]\label{proposition:1-1-2-image-properties}
像滤子确为 \(Y\) 上滤子，并且
\[
 f_*\mathcal F
 =\{B\subset Y:\text{存在 }F\in\mathcal F\text{ 使 }f(F)\subset B\}.
\]
若 \(g:Y\to Z\)，则 \((g\circ f)_*=g_*\circ f_*\)。
\end{proposition}
\begin{proof}
因为 \(f^{-1}(\varnothing)=\varnothing\notin\mathcal F\)，像滤子不含空集。逆像保持有限交与包含关系，
所以另两条滤子公理也成立。若 \(f^{-1}(B)\in\mathcal F\)，取 \(F=f^{-1}(B)\)，则 \(f(F)\subset B\)。
反向若 \(F\in\mathcal F\) 且 \(f(F)\subset B\)，则 \(F\subset f^{-1}(B)\)，向上封闭性给
\(f^{-1}(B)\in\mathcal F\)。最后
\((g\circ f)^{-1}(C)=f^{-1}(g^{-1}(C))\)，逐个 \(C\subset Z\) 比较定义即可得到函子性。
\end{proof}

\begin{proposition}[连续性的滤子判别]\label{proposition:1-1-2-3}
\(f:X\to Y\) 在 \(x\) 连续，当且仅当每个收敛到 \(x\) 的滤子 \(\mathcal F\) 都满足
\(f_*\mathcal F\to f(x)\)。
\end{proposition}
\begin{proof}
若 \(f\) 在 \(x\) 连续，给定 \(f(x)\) 的邻域 \(V\)，逆像 \(f^{-1}(V)\) 是 \(x\) 的邻域，故属于
\(\mathcal F\)。于是 \(V\in f_*\mathcal F\)。反过来只取 \(\mathcal F=\mathcal N(x)\)。若
\(f_*\mathcal N(x)\to f(x)\)，则对每个 \(f(x)\) 的邻域 \(V\)，有
\(f^{-1}(V)\in\mathcal N(x)\)，即 \(f^{-1}(V)\) 包含 \(x\) 的某个开邻域。这正是 \(f\) 在 \(x\)
连续。
\end{proof}

\begin{proposition}[聚点与细化]\label{proposition:1-1-2-4}
点 \(x\) 是滤子 \(\mathcal F\) 的聚点，当且仅当存在细化滤子
\(\mathcal G\supset\mathcal F\) 收敛到 \(x\)。
\end{proposition}
\begin{proof}
若 \(x\) 是聚点，集合族
\[
 \mathcal B=\{F\cap U:F\in\mathcal F,\ U\in\mathcal N(x)\}
\]
不含空集。两个基元素的交包含基元素
\((F_1\cap F_2)\cap(U_1\cap U_2)\)，故 \(\mathcal B\) 是滤子基。它生成的滤子 \(\mathcal G\)
同时包含 \(\mathcal F\) 和 \(\mathcal N(x)\)，所以 \(\mathcal G\to x\)。

反过来，若 \(\mathcal G\supset\mathcal F\) 且 \(\mathcal G\to x\)，则对 \(F\in\mathcal F\) 与
\(U\in\mathcal N(x)\)，有 \(F,U\in\mathcal G\)，所以 \(F\cap U\in\mathcal G\)。真滤子不含空集，
故 \(F\cap U\ne\varnothing\)，于是 \(x\) 是 \(\mathcal F\) 的聚点。
\end{proof}

\begin{example}[余有限滤子]\label{ex:1-1-2-1}
无限集合 \(X\) 上所有余有限集构成滤子。对通常序列 \(n\mapsto x_n\)，其尾滤子等于 \(\N\) 上余有限
滤子的像滤子；所以“\(n\to\infty\)”可以理解为一种大集合规则。
\end{example}

\begin{example}[邻域滤子]\label{ex:1-1-2-2}
\(\mathcal N(x)\) 是所有包含某个含 \(x\) 开集的子集所成集合族，而不只是开邻域本身。这样它才向上封闭。
上文已经逐项核对了滤子公理；连续性可写为
\(f_*\mathcal N(x)\supset\mathcal N(f(x))\)。
\end{example}

\begin{example}[概率一事件与依概率收敛]\label{ex:1-1-2-3}
本例只作第~5~章术语预告：概率空间、随机变量、依概率收敛及下面用到的概率运算都将在那里定义和证明，
本章后续不使用本例的结论。
在概率空间上，可测的概率一事件在可测集代数中构成滤子；把它们在 \(\mathcal P(\Omega)\) 中向上封闭，
可得到点集滤子。这表达“忽略零概率例外”的规则。若随机变量 \(X_n\) 依概率收敛到 \(X\)，则对每个
\(\varepsilon>0\)，数值 \(\Prob(\abs{X_n-X}<\varepsilon)\) 趋于 \(1\)；随 \(n\) 变化的这些事件本身
不因此组成滤子。两者的区别在于：概率一事件滤子固定在样本空间上，而依概率收敛控制的是一列事件的概率。
\end{example}

\begin{figure}[htbp]
\centering
\begin{tikzpicture}[node distance=1.45cm and 1.8cm,>=Stealth,every node/.style={align=center}]
 \node[book node] (n) {网\\带参数};
 \node[book node,right=of n] (f) {尾滤子\\最终大集合};
 \node[book node,right=of f] (c) {收敛与聚点\\邻域检验};
 \node[book node,below=of f] (p) {\((F,x)\) 成对索引\\精确恢复};
 \draw[book arrow] (n)--node[above]{忘掉参数}(f);
 \draw[book arrow] (f)--(c);
 \draw[book arrow] (f)--(p);
 \draw[book arrow] (p)--node[below]{\textsf{ZF}}(n);
\end{tikzpicture}
\caption{网与滤子保留相同的最终行为，但滤子忘掉参数化。}\label{fig:1-1-2-net-filter-translation}
\end{figure}

滤子的包含方向容易造成误会：成员越多，条件越强。像滤子不能只写成
\(\{f(F):F\in\mathcal F\}\)，因为这个像集族未必向上封闭。另一方面，给定 \(Y\) 上滤子
\(\mathcal G\)，逆像族 \(\{f^{-1}(G):G\in\mathcal G\}\) 可能含空集；它成为滤子基的充要条件是每个
\(G\) 都与 \(f(X)\) 相交。

\begin{exercise}
证明主滤子 \(\{A\subset X:E\subset A\}\) 为真滤子当且仅当 \(E\ne\varnothing\)。
\end{exercise}
\begin{exercise}
证明逆像族 \(\{f^{-1}(G):G\in\mathcal G\}\) 为滤子基当且仅当每个 \(G\in\mathcal G\) 与
\(f(X)\) 相交，并用常值映射构造失败例。
\end{exercise}
\begin{exercise}
证明滤子在 Hausdorff 空间中至多收敛到一点。
\end{exercise}
\begin{exercise}
用 \((F,x)\) 成对索引从给定滤子构造尾滤子恰等于它的网，并与逐 \(F\) 预选 \(x_F\in F\) 的写法比较。
\end{exercise}

滤子已经把“最终落入”改写为集合族的包含关系；要从一个不收敛的对象中提取稳定部分，还需说明何种
细化既保留全部尾部信息，又足以在每个集合及其补集之间作出选择。

\subsection{子网、超滤子与选择原则}\label{subsec:1-1-3}

从近似族中提取稳定部分，是分析中反复出现的动作。序列使用子序列；一般网需要一种既能越过每个旧阈值、
又不会不断返回早期指标的重参数化。滤子的对应动作是细化。若进一步要求每个集合 \(A\) 与其补集之间必须
作出一致选择，就得到超滤子。

\begin{remark}
子网在文献中有若干等价或近似的定义。这里采用最终共尾映射，因为它保证最终性质传递，并保证聚点能够产生
收敛子网。超滤子引理弱于完整选择公理这一集合论事实将在 \S\ref{subsec:1-4-3} 中准确表述；本小节只在
需要扩张滤子时明确说明所采用的原则。
\end{remark}

\begin{definition}[子网]\label{def:1-1-3-1}
网 \((y_\beta)_{\beta\in J}\) 称为 \((x_\alpha)_{\alpha\in I}\) 的\emph{子网}，若存在映射
\(\phi:J\to I\)，使 \(y_\beta=x_{\phi(\beta)}\)，且 \(\phi\) \emph{最终共尾}：对每个
\(\alpha_0\in I\)，存在 \(\beta_0\in J\)，使
\[
 \beta\ge\beta_0\quad\Longrightarrow\quad\phi(\beta)\ge\alpha_0.
\]
\end{definition}

\begin{example}[仅有共尾像还不够]\label{ex:1-1-3-3}
令 \(x_n=1/n\)，并定义 \(\phi(2k)=k\)、\(\phi(2k-1)=1\)。像集 \(\phi(\N)\) 在 \(\N\) 中共尾，
但 \(\phi\) 不是最终共尾：每个奇数指标仍映到 \(1\)。重参数后的网在奇数项恒取 \(1\)，因而不趋于 \(0\)。
\end{example}

\begin{proposition}[子网继承最终性质]\label{proposition:1-1-3-subnet-eventual}
若 \(P(x_\alpha)\) 最终成立，且 \((x_{\phi(\beta)})\) 是子网，则
\(P(x_{\phi(\beta)})\) 最终成立。特别地，收敛网的每个子网收敛到同一给定极限。
\end{proposition}
\begin{proof}
取 \(\alpha_0\) 使 \(\alpha\ge\alpha_0\) 时 \(P(x_\alpha)\) 成立。最终共尾性给出 \(\beta_0\)，使
\(\beta\ge\beta_0\) 时 \(\phi(\beta)\ge\alpha_0\)。于是这些 \(\beta\) 都满足
\(P(x_{\phi(\beta)})\)。对收敛网，把 \(P\) 取为“网值属于给定极限的邻域”即可。
\end{proof}

\begin{definition}[超滤子]\label{def:1-1-3-2}
真滤子 \(\mathcal U\) 若在包含关系下极大，称为\emph{超滤子}。
\end{definition}

\begin{definition}[超滤子极限]\label{def:1-1-3-3}
若超滤子 \(\mathcal U\) 收敛到 \(x\)，称 \(x\) 为其极限。
\end{definition}

\begin{proposition}[超滤子二分刻画]\label{proposition:1-1-3-2}
真滤子 \(\mathcal U\) 为超滤子，当且仅当对每个 \(A\subset X\)，恰有 \(A\in\mathcal U\) 或
\(A^c\in\mathcal U\) 之一成立。
\end{proposition}
\begin{proof}
先设 \(\mathcal U\) 极大。若 \(A\notin\mathcal U\) 且对每个 \(U\in\mathcal U\) 都有
\(A\cap U\ne\varnothing\)，则
\[
  \mathcal B_A=\{A\cap U:U\in\mathcal U\}
\]
不含空集，并且
\[
  (A\cap U_1)\cap(A\cap U_2)=A\cap(U_1\cap U_2)\in\mathcal B_A.
\]
故 \(\mathcal B_A\) 是滤子基；它生成的真滤子同时包含 \(\mathcal U\) 与 \(A\)，从而严格细化
\(\mathcal U\)，与极大性矛盾。因此存在 \(U\in\mathcal U\) 使 \(A\cap U=\varnothing\)。于是
\(U\subset A^c\)，向上封闭性给出 \(A^c\in\mathcal U\)。二者不能同时属于 \(\mathcal U\)，因为
\[
  A,A^c\in\mathcal U\quad\Longrightarrow\quad
  \varnothing=A\cap A^c\in\mathcal U.
\]

反过来，设二分性质成立，且真滤子 \(\mathcal G\supset\mathcal U\)。若存在
\(A\in\mathcal G\setminus\mathcal U\)，则 \(A^c\in\mathcal U\subset\mathcal G\)，于是
\(\varnothing=A\cap A^c\in\mathcal G\)，矛盾。因此 \(\mathcal G=\mathcal U\)，即 \(\mathcal U\) 极大。
\end{proof}

\begin{remark}
超滤子的聚点自动成为极限。若 \(x\) 是 \(\mathcal U\) 的聚点而 \(U\in\mathcal N(x)\) 不属于
\(\mathcal U\)，二分性质给 \(U^c\in\mathcal U\)，但 \(U^c\cap U=\varnothing\)，与聚点定义矛盾。
\end{remark}

\begin{example}[自由超滤子]\label{ex:1-1-3-1}
无限集合 \(X\) 上的余有限滤子若扩张为超滤子 \(\mathcal U\)，则 \(\mathcal U\) 不含有限集。事实上，
若有限集 \(F\in\mathcal U\)，则余有限集 \(X\setminus F\) 也属于 \(\mathcal U\)，从而
\[
  \varnothing=F\cap(X\setminus F)\in\mathcal U,
\]
矛盾。这样的 \(\mathcal U\) 称为自由超滤子。它的存在来自超滤子引理；这里不把某个自由超滤子当作
可显式列出的对象。
\end{example}

\begin{example}[有限划分与可数划分]\label{ex:1-1-3-2}
任一 \(\N\) 上超滤子恰包含偶数集或奇数集之一。更一般地，若
\(X=A_1\sqcup\cdots\sqcup A_m\)，且没有 \(A_j\) 属于 \(\mathcal U\)，二分性质便给出
\(A_j^c\in\mathcal U\)（\(1\le j\le m\)），从而
\[
  \varnothing=\bigcap_{j=1}^m A_j^c\in\mathcal U,
\]
矛盾；而两个不同胞腔不能同时属于 \(\mathcal U\)，因为它们的交为空。因此有限划分中恰有一个胞腔
属于超滤子。有限性不能删除：自由超滤子包含 \(\N\)，却不包含任何单点，因而不选择可数单点划分中的
胞腔。
\end{example}

\begin{theorem}[聚点--子网定理]\label{theorem:1-1-3-1}
点 \(x\) 是网 \((x_\alpha)_{\alpha\in I}\) 的聚点，当且仅当该网有一个收敛到 \(x\) 的子网。
\end{theorem}
\begin{proof}
若有子网 \(x_{\phi(\beta)}\to x\)，给定 \(x\) 的邻域 \(U\) 和旧阈值 \(\alpha_0\)，子网收敛给出
\(\beta_1\)，使 \(\beta\ge\beta_1\) 时 \(x_{\phi(\beta)}\in U\)；最终共尾性给出 \(\beta_2\)，使
\(\beta\ge\beta_2\) 时 \(\phi(\beta)\ge\alpha_0\)。取两者的共同后继 \(\beta\)，便得到一个
\(\alpha=\phi(\beta)\ge\alpha_0\) 且 \(x_\alpha\in U\)。故 \(x\) 是聚点。

反过来设 \(x\) 是聚点。令
\[
 J=\{(\alpha,U):\alpha\in I,\ U\in\mathcal N(x),\ x_\alpha\in U\},
\]
并定义
\[
 (\alpha,U)\preceq(\beta,V)
 \quad\Longleftrightarrow\quad
 \alpha\le\beta\text{ 且 }V\subset U.
\]
聚点定义保证 \(J\ne\varnothing\)。给定 \((\alpha,U),(\beta,V)\)，先取旧指标共同后继 \(\delta\)。
因为网频繁进入 \(U\cap V\)，存在 \(\gamma\ge\delta\) 使 \(x_\gamma\in U\cap V\)。于是
\((\gamma,U\cap V)\) 是共同后继，故 \(J\) 有向。

令 \(\phi(\alpha,U)=\alpha\)。给定 \(\alpha_0\)，取任意邻域 \(W\)；聚点性给出
\(\alpha_1\ge\alpha_0\) 且 \(x_{\alpha_1}\in W\)，所以 \((\alpha_1,W)\in J\)。若
\((\beta,V)\succeq(\alpha_1,W)\)，则 \(\phi(\beta,V)=\beta\ge\alpha_1\ge\alpha_0\)。故
\(\phi\) 最终共尾。

子网值为 \(x_{\phi(\alpha,U)}=x_\alpha\)。给定邻域 \(W\ni x\)，取 \((\alpha_1,W)\in J\) 如上；
在该指标之后，邻域坐标 \(V\subset W\)，且子网值 \(x_\beta\in V\subset W\)。故子网收敛到 \(x\)。
整个构造使用所有可行对，属于 \textsf{ZF} 证明。
\end{proof}

下面的证明使用一条外部集合论原则。Zorn 引理断言：若一个非空偏序集的每条链都有上界，则该偏序集
有极大元。在 \textsf{ZF} 上，Zorn 引理与选择公理 \textsf{AC} 等价；本书不证明这项集合论等价，
只在下述证明中明确调用 Zorn 引理。超滤子引理本身可以作为较弱的 \textsf{UFL} 原则单独采用。

\begin{theorem}[超滤子引理]\label{theorem:1-1-3-3}
每个真滤子都包含于某个超滤子。
\end{theorem}
\begin{proof}
在通常采用选择公理的语境中，可由 Zorn 引理证明。固定真滤子 \(\mathcal F\)，令 \(\mathscr P\) 为所有
包含 \(\mathcal F\) 的真滤子，按包含排序。若 \(\mathscr C\subset\mathscr P\) 是链，令
\(\mathcal H=\bigcup_{\mathcal G\in\mathscr C}\mathcal G\)。空集不属于 \(\mathcal H\)。若
\(A,B\in\mathcal H\)，存在链中两个滤子分别包含它们；链的线性次序使其中一个滤子同时包含 \(A,B\)，
所以 \(A\cap B\in\mathcal H\)。向上封闭性也由同一成员滤子继承。因此 \(\mathcal H\) 是链的上界。
Zorn 引理给出极大成员，它就是包含 \(\mathcal F\) 的超滤子。

这份证明使用 \textsf{AC}。命题本身通常称超滤子引理，并可作为较弱原则 \textsf{UFL} 采用；上述证明
不说明命题必须使用完整选择公理。
\end{proof}

\begin{proposition}[超滤子像]\label{proposition:1-1-3-4}
若 \(f:X\to Y\)，\(\mathcal U\) 是 \(X\) 上超滤子，则 \(f_*\mathcal U\) 是 \(Y\) 上超滤子。
\end{proposition}
\begin{proof}
像滤子是真滤子。对任意 \(B\subset Y\)，超滤子的二分性质应用于 \(f^{-1}(B)\) 给出
\(f^{-1}(B)\in\mathcal U\) 或
\(X\setminus f^{-1}(B)=f^{-1}(B^c)\in\mathcal U\)，且二者不能同时成立。依像滤子定义，
\(B,B^c\) 中恰有一个属于 \(f_*\mathcal U\)。由二分刻画，像滤子为超滤子。
\end{proof}

\begin{figure}[htbp]
\centering
\begin{tikzpicture}[node distance=1.45cm and 1.8cm,>=Stealth,every node/.style={align=center}]
 \node[book node] (a) {网的聚点};
 \node[book node,right=of a] (b) {收敛子网};
 \node[book node,below=of a] (c) {尾滤子的聚点};
 \node[book node,right=of c] (d) {收敛细化};
 \node[book node,below right=of c] (e) {超滤子扩张};
 \draw[book arrow,<->] (a)--node[above]{\textsf{ZF}}(b);
 \draw[book arrow,<->] (a)--node[left]{翻译}(c);
 \draw[book arrow,<->] (c)--node[above]{\textsf{ZF}}(d);
 \draw[book dashed arrow] (d)--node[right]{\textsf{UFL}}(e);
\end{tikzpicture}
\caption{稳定部分的三种语言以及极大化发生的位置。}\label{fig:1-1-3-cluster-subnet-ultrafilter}
\end{figure}

子网不必由原指标集的子集得到；关键是最终共尾映射。超滤子极限也不自动唯一，唯一性仍需 Hausdorff 条件。
超滤子的二分只控制有限划分，这一点将在测度论的可数可加性出现时形成重要对照。

\begin{exercise}
证明收敛网的每个子网收敛到同一给定极限。
\end{exercise}
\begin{exercise}
证明对超滤子 \(\mathcal U\)，有限并 \(A_1\cup\cdots\cup A_n\in\mathcal U\) 当且仅当某个
\(A_j\in\mathcal U\)。
\end{exercise}
\begin{exercise}
从网的聚点出发，重新构造聚点--子网定理中的有向集，并逐项检查有向性、最终共尾性与收敛性。
\end{exercise}
\begin{exercise}
设 \(X\) 第一可数。对普通序列证明：\(x\) 是序列聚点，当且仅当存在收敛到 \(x\) 的子序列。
\end{exercise}

下一节把紧性表述为：所有有向观测都能产生聚点，或者所有超滤子都能找到极限。
